\section{结论}
深度学习正在成为智慧能源研究中的关键支撑技术，尤其在光伏和风电等新能源预测任务中表现出明显优势。通过对代表性文献的分析可以看出，LSTM、GRU、CNN-LSTM、TCN 等模型能够较好地处理强时序依赖与复杂非线性关系，在短时预测场景中具有较高应用价值。同时，领域知识融合、元学习和强化学习调度案例表明，深度学习的价值并不止于提高单一预测指标，而是能够进一步服务于储能控制、微电网运行和能源管理优化。

但也应看到，深度学习并不是智慧能源问题的“万能解”。如果数据质量不足、模型缺乏约束、结果难以解释，那么即使测试集误差较小，也未必能在真实系统中稳定落地。因此，本文最终的判断是：深度学习在智慧能源中的核心意义，不是把模型做得更复杂，而是利用数据驱动方法提升系统对不确定性的感知、预测和应对能力。未来真正高价值的方向，将是可解释、可迁移、可部署，并能直接支撑能源决策的深度学习方法。
